\documentclass[12pt]{article}
\usepackage{amsmath, amssymb, amsthm}
\usepackage{geometry}
\usepackage{hyperref}
\usepackage{graphicx}
\usepackage{bm}
\usepackage{xcolor}
\usepackage{mathtools}
\usepackage{physics}
\geometry{margin=1in}

\newcommand{\R}{\mathbb{R}}
\newcommand{\Z}{\mathbb{Z}}
\newcommand{\N}{\mathbb{N}}
\newcommand{\eps}{\varepsilon}
\newcommand{\half}{\tfrac{1}{2}}
\newcommand{\pderiv}[2]{\frac{\partial #1}{\partial #2}}
\newcommand{\avec}{\mathbf{a}}
\newcommand{\fvec}{\mathbf{f}}
\newcommand{\uvec}{\hat{\mathbf{u}}}
\newcommand{\vvec}{\hat{\mathbf{v}}}
\newcommand{\nvec}{\hat{\mathbf{n}}}
\newcommand{\xvec}{\mathbf{x}}
\DeclareMathOperator{\sgn}{sgn}
\DeclareMathOperator{\atantwo}{atan2}

\title{\textbf{Mathematical Notes: Rounded-Polygon Packing Model}\\[0.5em]
\large Derivations of all algorithms in \texttt{pyCudaPolygon}}
\author{}
\date{\today}

\begin{document}
\maketitle
\tableofcontents
\newpage

%% ============================================================
\section{Overview}
%% ============================================================

This document gives step-by-step derivations of every piece of
mathematics used in the \texttt{pyCudaPolygon} CUDA code for
soft rounded-polygon packing.  The chain is:

\begin{enumerate}
  \item \textbf{Polygon geometry}: equilateral closed polygons, shape index.
  \item \textbf{Rounded-polygon boundary}: inward arc rounding at each vertex.
  \item \textbf{Arc geometry} (\texttt{computeArcGeom}): tangent points,
        arc centre, angle span.
  \item \textbf{Intersection area as energy}: Green's theorem applied to arcs
        and straight edges.
  \item \textbf{Force as gradient of intersection area}: analytic gradient
        via the arc-geometry chain rule (\texttt{arcGeomGrad}).
  \item \textbf{SHAKE}: simultaneous Newton-step constraint projection for
        equilateral edge lengths and area.
  \item \textbf{Force projection}: SVD + Gram–Schmidt to remove all
        constraint-space force components.
  \item \textbf{FIRE minimization}: velocity-Verlet FIRE with SHAKE rederivation.
  \item \textbf{Shape generation}: SLSQP for random equilateral polygons with
        prescribed shape index.
  \item \textbf{Bidisperse packing setup}: area formulas from packing fraction
        and perimeter ratio.
\end{enumerate}

Throughout, positions live in a periodic $[0,1)^2$ box with minimum-image convention
\[
  \mathrm{wrap}(x) = x - \lfloor x + \tfrac{1}{2} \rfloor.
\]

%% ============================================================
\section{Polygon Geometry}
%% ============================================================

\subsection{Backbone polygon}

A polygon $P$ with $N$ vertices is represented by its backbone vertices
$\mathbf{a}_0, \mathbf{a}_1, \ldots, \mathbf{a}_{N-1} \in [0,1)^2$ (periodic box),
traversed counterclockwise.  Edges are labelled $E_k$, running from $\mathbf{a}_k$
to $\mathbf{a}_{k+1}$ (indices mod $N$).

\subsubsection{Edge vectors}

The (wrapped) edge vectors are
\[
  \mathbf{u}_k = \mathrm{wrap}(\mathbf{a}_k - \mathbf{a}_{k-1}), \qquad
  \mathbf{v}_k = \mathrm{wrap}(\mathbf{a}_{k+1} - \mathbf{a}_k).
\]
Unit vectors: $\uvec_k = \mathbf{u}_k / |\mathbf{u}_k|$, \;
$\vvec_k = \mathbf{v}_k / |\mathbf{v}_k|$.

\subsubsection{Area (shoelace formula)}

The signed area of the backbone polygon (positive for CCW) is
\[
  A_{\mathrm{bb}} = \frac{1}{2} \sum_{k=0}^{N-1}
      \bigl(x_k\, y_{k+1} - x_{k+1}\, y_k\bigr),
\]
where $(x_k, y_k) = \mathbf{a}_k$ and all coordinates are interpreted in a
single minimum-image frame anchored at $\mathbf{a}_0$.

\subsubsection{Equilateral constraint}

SHAKE enforces $|\mathbf{v}_k| = L_k = $ constant for all $k$.  When all $L_k = L$
the polygon is \emph{equilateral} (equal edge lengths but not necessarily equiangular).

\subsection{Shape index}

For a polygon with perimeter $P = \sum_k L_k$ and backbone area $A$, the
\emph{shape index} (also called the reduced perimeter) is
\[
  \kappa = \frac{P}{\sqrt{A}}.
\]
A circle has $\kappa = 2\sqrt{\pi} \approx 3.545$; a square has $\kappa = 4$;
elongated or irregular shapes have larger $\kappa$.  In the bidisperse packing,
we use $\kappa = 3.7$ for both polygon sizes.

\subsection{Exterior angle at a vertex}

At vertex $\mathbf{a}_k$ the incoming direction is $\uvec_k$ and the outgoing
direction is $\vvec_k$.  Define
\[
  c_k = \uvec_k \times \vvec_k = (\hat u_x)(\hat v_y) - (\hat u_y)(\hat v_x), \qquad
  d_k = \uvec_k \cdot \vvec_k = \hat u_x \hat v_x + \hat u_y \hat v_y.
\]
The \emph{exterior angle} (the angle you turn through at vertex $k$, positive for a
left turn in a CCW polygon) is
\[
  \varphi_k = \mathrm{atan2}(c_k, d_k).
\]
For a convex polygon traversed CCW, $\varphi_k > 0$ for all $k$ and
$\sum_k \varphi_k = 2\pi$.

%% ============================================================
\section{Rounded-Polygon Boundary}
%% ============================================================

\subsection{Inward arc rounding}

At each vertex $\mathbf{a}_k$ we replace the sharp corner with a circular arc of
radius $\delta > 0$.  The arc is tangent to both adjacent (original) edges and
curves \emph{inward} (toward the polygon interior).  The resulting boundary
alternates between straight edge sections and smooth arcs:

\[
  \partial \tilde P = K_0, E_0, K_1, E_1, \ldots, K_{N-1}, E_{N-1},
\]
where $K_k$ is the arc at vertex $\mathbf{a}_k$ and $E_k$ is the straight section
of edge $k$ between adjacent tangent points.

This is \textbf{not} the Minkowski sum $P \oplus D(\delta)$ (which rounds outward);
it is instead an \emph{inscribed} rounding that cuts each corner.

\subsection{Tangent points and arc centre}
\label{sec:arcgeom}

The following derivation corresponds exactly to \texttt{computeArcGeom} in
\texttt{kernels.cuh}.

\paragraph{Tangent set-back $\ell_k$.}
The arc at $\mathbf{a}_k$ must be tangent to the incoming edge (direction $\uvec_k$)
and to the outgoing edge (direction $\vvec_k$).  Let the tangent point on the incoming
edge be $\mathbf{a}_k^- = \mathbf{a}_k - \ell_k \uvec_k$ and on the outgoing edge
be $\mathbf{a}_k^+ = \mathbf{a}_k + \ell_k \vvec_k$.

For a circle of radius $\delta$ centred at $\mathbf{z}_k$ to be tangent to the
incoming edge at $\mathbf{a}_k^-$, we need:
\[
  (\mathbf{z}_k - \mathbf{a}_k^-) \perp \uvec_k \quad \text{and} \quad
  |\mathbf{z}_k - \mathbf{a}_k^-| = \delta.
\]
By symmetry (same tangency condition on outgoing edge), $|\mathbf{z}_k - \mathbf{a}_k^+|
= \delta$ as well.  The tangent-setback formula follows from the two perpendicularity
conditions:

Write $\mathbf{z}_k = \mathbf{a}_k + \alpha \uvec_k + \beta \nvec_k$ where
$\nvec_k$ is a unit vector perpendicular to $\uvec_k$.  For the arc to be tangent to
the incoming edge: the component of $(\mathbf{z}_k - \mathbf{a}_k^-)$ along $\uvec_k$
must vanish:
\[
  \alpha + \ell_k = 0 \implies \alpha = -\ell_k.
\]
For the arc to be tangent to the outgoing edge (direction $\vvec_k$): the component of
$(\mathbf{z}_k - \mathbf{a}_k^+)$ along $\vvec_k$ must vanish.  Writing
$\vvec_k = d_k \uvec_k + c_k \nvec_k$ (resolved in the $\uvec_k$, $\nvec_k$ frame):
\[
  (-\ell_k - \ell_k d_k) + c_k \beta = 0 \implies
  \beta = \frac{\ell_k (1 + d_k)}{c_k}.
\]
Using $|\mathbf{z}_k - \mathbf{a}_k^-| = \delta$:
\[
  \alpha^2 + \beta^2 - 2\alpha\ell_k + \ell_k^2 = \delta^2 \quad\Longrightarrow\quad
  \ell_k^2 + \beta^2 = \delta^2.
\]
Since $\beta = \ell_k(1+d_k)/c_k$ and $\beta^2 = \delta^2 - \ell_k^2$:
\[
  \ell_k^2 \left[1 + \frac{(1+d_k)^2}{c_k^2}\right] = \delta^2 \quad\Longrightarrow\quad
  \ell_k^2 \cdot \frac{c_k^2 + (1+d_k)^2}{c_k^2} = \delta^2.
\]
Expanding $c_k^2 + (1+d_k)^2 = 1 - d_k^2 + 1 + 2d_k + d_k^2 = 2(1 + d_k)$:
\[
  \ell_k^2 \cdot \frac{2(1+d_k)}{c_k^2} = \delta^2
  \quad\Longrightarrow\quad
  \boxed{\ell_k = \frac{\delta\, |c_k|}{1 + d_k} = \delta \tan\!\left(\frac{\varphi_k}{2}\right).}
\]
The last equality uses the half-angle identity $\tan(\varphi/2) = \sin\varphi/(1+\cos\varphi)$
together with $|c_k| = \sin\varphi_k$, $d_k = \cos\varphi_k$.

\paragraph{Validity condition.}
For the tangent points to lie within the edges (so adjacent arcs do not overlap),
we need $\ell_k + \ell_{k-1} < L_{k-1}$ for every edge $k-1$.  Since SHAKE maintains
equilateral edges, a sufficient condition is
\[
  2\,\delta\,\tan\!\left(\frac{\varphi_k}{2}\right) < L \quad \text{for all } k.
\]
For a regular $N$-gon this gives $\delta < L/(2\tan(\pi/N))$.

\paragraph{Arc centre.}
From $\alpha = -\ell_k$ and $\beta = \ell_k(1+d_k)/c_k = \delta|c_k|/c_k = \delta\,\sgn(c_k)$:
\[
  \mathbf{z}_k = \mathbf{a}_k - \ell_k \uvec_k + \delta\,\sgn(c_k)\, \nvec_k.
\]
Using $\nvec_k = (-\hat u_y, \hat u_x)$ (the 90$^\circ$ CCW rotation of $\uvec_k$) and
noting that for a CCW polygon $c_k > 0$, this simplifies to $\beta = \delta$, i.e.:
\[
  \mathbf{z}_k = \mathbf{a}_k - \ell_k \uvec_k + \delta\nvec_k.
\]
This can be rewritten as
\[
  \boxed{\mathbf{z}_k = \mathbf{a}_k + \frac{\delta(\vvec_k - \uvec_k)}{|c_k|},}
\]
which is what the code computes.  \emph{Proof}: $-\ell_k \uvec_k + \delta \nvec_k
= -\frac{\delta|c_k|}{1+d_k}\uvec_k + \delta(-\hat u_y, \hat u_x)$.
Meanwhile $(\vvec_k - \uvec_k)/|c_k|$: write $\vvec_k - \uvec_k = (d_k-1)\uvec_k + c_k\nvec_k$,
so $(\vvec_k-\uvec_k)/|c_k| = (d_k-1)\uvec_k/|c_k| + \sgn(c_k)\nvec_k$.
Multiplying by $\delta$: $\delta(d_k-1)/|c_k| \cdot \uvec_k + \delta\sgn(c_k)\nvec_k$.
Since $(d_k-1)/|c_k| = -(1-d_k)/(|c_k|) = -(1+d_k-2)/(|c_k|)$\ldots a cleaner route:
with $\varphi = \varphi_k$, $c = \sin\varphi$, $d = \cos\varphi$:
\begin{align*}
  \frac{\delta(\vvec-\uvec)}{|c|} &= \frac{\delta}{|\sin\varphi|}
    \bigl[(\cos\varphi - 1)\uvec + \sin\varphi\,\nvec\bigr] \\
  &= \delta\,\frac{\cos\varphi-1}{\sin\varphi}\,\uvec + \delta\,\nvec
   = -\delta\tan\!\tfrac{\varphi}{2}\,\uvec + \delta\,\nvec
   = -\ell_k\uvec_k + \delta\nvec_k.
\end{align*}

\paragraph{Arc angle span.}
The arc runs from $\mathbf{a}_k^-$ to $\mathbf{a}_k^+$ centred at $\mathbf{z}_k$.
The starting angle (angle of $\mathbf{a}_k^- - \mathbf{z}_k$ measured from the
positive $x$-axis) is
\[
  \phi_{0,k} = \mathrm{atan2}\!\bigl(a^-_{k,y} - z_{k,y},\; a^-_{k,x} - z_{k,x}\bigr).
\]
The arc sweeps through $\Delta\varphi_k = \varphi_k$ (the exterior angle), i.e.\
the arc spans exactly the same angle as the turn at the vertex.

\subsection{Straight edge section}

The edge section $E_k$ (the flat part of edge $k$) runs from $\mathbf{a}_k^+$ to
$\mathbf{a}_{k+1}^-$, a straight line segment of length $L_k - \ell_k - \ell_{k+1}$.

\subsection{Area of rounded polygon}

The area of the inscribed rounded polygon $\tilde P$ is
\[
  A_{\tilde P} = A_{\mathrm{bb}} - \sum_{k=0}^{N-1} \Delta A_k,
\]
where $\Delta A_k$ is the area of the corner region removed at vertex $k$.  Each
removed corner is the region bounded by the arc $K_k$ and the two straight segments
from $\mathbf{a}_k$ to $\mathbf{a}_k^\pm$.  Its area is

\[
  \Delta A_k = \frac{1}{2}\ell_k^2 \tan\!\left(\frac{\varphi_k}{2}\right)
              - \frac{1}{2}\delta^2(\varphi_k - \sin\varphi_k).
\]
\emph{Derivation}: The two triangles from $\mathbf{a}_k$ to the tangent points
have total area $\frac{1}{2}\ell_k^2 \sin\varphi_k / \cos^2(\varphi_k/2)
= \frac{1}{2}\ell_k^2\tan(\varphi_k/2)$ (using the tangency geometry).
The circular sector subtracted is $\frac{1}{2}\delta^2\varphi_k$ and the triangles
from $\mathbf{z}_k$ to $\mathbf{a}_k^\pm$ together contribute $\frac{1}{2}\delta^2\sin\varphi_k$.
The net result is the formula above.

%% ============================================================
\section{Intersection Area as the Energy}
%% ============================================================

\subsection{Setup}

For two rounded polygons $\tilde A$ and $\tilde B$ (corresponding to backbone polygons
$A$ and $B$), the potential energy is
\[
  E_{AB} = \mathrm{Area}(\tilde A \cap \tilde B).
\]
The total energy is $E = \sum_{A < B} E_{AB}$ summed over interacting pairs.

\subsection{Green's theorem}

Green's theorem (the ``shoelace'' form) states that for a region $\Omega$ with
boundary $\partial\Omega$ traversed counterclockwise,
\[
  \mathrm{Area}(\Omega) = \frac{1}{2}\oint_{\partial\Omega}
    \bigl(x\,\dd y - y\,\dd x\bigr).
\]
We compute $\mathrm{Area}(\tilde A \cap \tilde B)$ by integrating over the boundary
of $\tilde A \cap \tilde B$.  The intersection boundary consists of alternating arcs
and straight segments from $\partial\tilde A$ and $\partial\tilde B$.

\paragraph{Key decomposition.}
Rather than assembling the full intersection boundary (which requires polygon clipping),
the CUDA kernel decomposes the computation per feature of $\tilde A$:

\begin{quote}
For each feature $F$ of $\partial\tilde A$ (either arc $K_i$ or edge $E_i$):
\begin{enumerate}
  \item Clip $F$ against $\tilde B$: find the sub-intervals where $F$ lies \emph{inside}
        $\tilde B$.
  \item For each such sub-interval $[t_0, t_1]$, add the corresponding signed area
        contribution $\delta\!\mathcal{A}$ to the energy.
\end{enumerate}
\end{quote}

This is valid because, when we apply Green's theorem to $\tilde A \cap \tilde B$,
the only contributions come from the parts of the boundary that lie inside the other
shape.  Summing over all features of $\tilde A$ gives the full intersection area.
(The contributions from $\partial\tilde B$ inside $\tilde A$ are captured when the
roles of $A$ and $B$ are reversed; both halves are summed in the final accumulation
via \texttt{atomicAdd}.)

\subsection{Contribution from a straight edge segment}
\label{sec:edge_contrib}

Edge section $E_i$ runs from $\mathbf{a}_i^+ \equiv \mathbf{P}$ to
$\mathbf{a}_{i+1}^- \equiv \mathbf{Q}$ parametrically as
\[
  \boldsymbol{\gamma}(t) = (1-t)\mathbf{P} + t\mathbf{Q}, \quad t \in [0,1].
\]
Suppose the sub-interval $[t_0, t_1]$ lies inside $\tilde B$, with endpoints
$\mathbf{X}_0 = \boldsymbol{\gamma}(t_0)$ and $\mathbf{X}_1 = \boldsymbol{\gamma}(t_1)$.
The Green's theorem integrand over this segment is
\[
  \dd\mathcal{A} = \frac{1}{2}(x\,\dd y - y\,\dd x).
\]
On the straight segment:
\[
  x = (1-t)P_x + tQ_x, \quad \dd x = (Q_x - P_x)\,\dd t,
\]
and similarly for $y$.  Therefore:
\[
  x\,\dd y - y\,\dd x = \bigl[(1-t)P_x + tQ_x\bigr](Q_y - P_y)\,\dd t
                       - \bigl[(1-t)P_y + tQ_y\bigr](Q_x - P_x)\,\dd t.
\]
Integrating from $t_0$ to $t_1$:
\begin{align*}
  \delta\mathcal{A}_{\mathrm{edge}} &= \frac{1}{2}\int_{t_0}^{t_1}
    \bigl[(1-t)(P_x Q_y - P_y Q_x) + t(Q_x P_y - Q_y P_x) \cdot (-1)
    + \ldots\bigr]\,\dd t.
\end{align*}
After simplification, the result depends only on the \emph{endpoint coordinates}
$\mathbf{X}_0$ and $\mathbf{X}_1$ of the inside sub-interval:
\[
  \boxed{
  \delta\mathcal{A}_{\mathrm{edge}}([\mathbf{X}_0, \mathbf{X}_1])
  = \frac{1}{2}\bigl(X_{0x}\,X_{1y} - X_{0y}\,X_{1x}\bigr).
  }
\]
\emph{Proof}: For a straight segment from $\mathbf{X}_0$ to $\mathbf{X}_1$:
$x = (1-s)X_{0x} + s X_{1x}$, $\dd x = (X_{1x} - X_{0x})\,\dd s$.  Then:
\begin{align*}
  \int_0^1 (x\,\dd y - y\,\dd x) &=
  \int_0^1 \bigl[(1-s)X_{0x} + sX_{1x}\bigr](X_{1y}-X_{0y})\,\dd s \\
  &\quad - \int_0^1\bigl[(1-s)X_{0y}+sX_{1y}\bigr](X_{1x}-X_{0x})\,\dd s \\
  &= \tfrac{1}{2}(X_{0x}+X_{1x})(X_{1y}-X_{0y})
    - \tfrac{1}{2}(X_{0y}+X_{1y})(X_{1x}-X_{0x}) \\
  &= X_{0x}X_{1y} - X_{0y}X_{1x}.
\end{align*}
This matches the kernel line: \texttt{localEnergy += 0.5*(X0x*X1y - X0y*X1x)}.

\subsection{Contribution from an arc segment}
\label{sec:arc_contrib}

Arc $K_i$ at vertex $\mathbf{a}_i$ is a circle of radius $\delta$ centred at
$\mathbf{z}_i$, parametrised by
\[
  \boldsymbol{\gamma}(\tau) = \mathbf{z}_i + \delta\bigl(\cos(\phi_{0,i} + \tau\Delta\varphi_i),\;
                                         \sin(\phi_{0,i} + \tau\Delta\varphi_i)\bigr),
  \quad \tau \in [0,1].
\]
Suppose the inside sub-interval is $[\tau_0, \tau_1]$ with endpoints $\mathbf{X}_0$,
$\mathbf{X}_1$ and spanning angle $\Delta\psi = (\tau_1 - \tau_0)\Delta\varphi_i$.

The Green's integrand on a circle of radius $\delta$ centred at $\mathbf{z}$:
\[
  x\,\dd y - y\,\dd x = (z_x + \delta\cos\phi)(-\delta\sin\phi\,\dd\phi)
    - (z_y + \delta\sin\phi)(\delta\cos\phi\,\dd\phi)
  = -\delta\bigl(z_x\sin\phi + z_y\cos\phi + \delta\bigr)\,\dd\phi.
\]
Integrate from $\phi_0 = \phi_{0,i} + \tau_0\Delta\varphi_i$ to
$\phi_1 = \phi_{0,i} + \tau_1\Delta\varphi_i$:
\[
  \int (x\,\dd y - y\,\dd x)
    = \delta\bigl[z_x\cos\phi - z_y\sin\phi\bigr]_{\phi_0}^{\phi_1}
    - \delta^2(\phi_1 - \phi_0).
\]
Using $\delta\cos\phi_{0/1} = X_{0/1,x} - z_x$ and $\delta\sin\phi_{0/1} = X_{0/1,y} - z_y$:
\begin{align*}
  \int (x\,\dd y - y\,\dd x)
  &= z_x(X_{1x}-z_x) - z_y(X_{1y}-z_y) - z_x(X_{0x}-z_x) + z_y(X_{0y}-z_y)
     - \delta^2\Delta\psi \\
  &= z_x(X_{1x} - X_{0x}) - z_y(X_{1y} - X_{0y}) - \delta^2\Delta\psi.
\end{align*}
We can write $z_x X_{1x} - z_x X_{0x} = X_{0x}X_{1y} - X_{0y}X_{1x}
+ z_x(X_{1x}-X_{0x}) - z_y(X_{1y}-X_{0y}) + \ldots$ — it's cleaner to
directly compute.  The shoelace contribution for the chord plus the circular cap:
\[
  \mathrm{Area}(\text{arc sector}) = \frac{1}{2}(X_{0x}X_{1y} - X_{0y}X_{1x})
  + \frac{1}{2}\delta^2(\Delta\psi - \sin\Delta\psi).
\]
\emph{Derivation}: The signed area swept by the arc from $\mathbf{X}_0$ to $\mathbf{X}_1$
centred at $\mathbf{z}$ can be split into (i) the triangle $\mathbf{0}$–$\mathbf{X}_0$–$\mathbf{X}_1$
with area $\frac{1}{2}(X_{0x}X_{1y}-X_{0y}X_{1x})$ and (ii) the circular
segment (area enclosed between the chord $\mathbf{X}_0\mathbf{X}_1$ and the arc).
The circular segment area is $\frac{1}{2}\delta^2(\Delta\psi - \sin\Delta\psi)$.
Combining:
\[
  \boxed{
  \delta\mathcal{A}_{\mathrm{arc}} = \frac{1}{2}\bigl(X_{0x}X_{1y} - X_{0y}X_{1x}\bigr)
  + \frac{1}{2}\delta^2(\Delta\psi - \sin\Delta\psi).
  }
\]
This matches the kernel lines:
\begin{verbatim}
  localEnergy += 0.5*(X0x*X1y - X0y*X1x)
               + 0.5*delta*delta*(dpsi - sin(dpsi));
\end{verbatim}

\subsection{Point-in-smoothed-polygon test}

To determine whether a midpoint of a feature segment lies inside $\tilde B$, the
code uses a horizontal ray-casting algorithm (\texttt{pipSmoothed}).  A horizontal
ray from point $(p_x, p_y)$ toward $+x$ is cast, and the number of crossings with
$\partial\tilde B$ (both arc and straight segments) is counted.  An odd count means
the point is inside.

For a straight edge from $\mathbf{b}^+_j$ to $\mathbf{b}^-_{j+1}$: standard
ray–segment crossing test.

For an arc at vertex $j$ with centre $(z_x, z_y)$, radius $\delta$: the ray
$y = p_y$ intersects the circle at $x = z_x \pm \sqrt{\delta^2 - (p_y - z_y)^2}$
(if the discriminant is non-negative).  Each intersection is checked to lie on the
arc (i.e., the corresponding angle $\phi = \mathrm{atan2}(p_y - z_y, x - z_x)$ falls
in the arc range $[\phi_0, \phi_0 + \Delta\varphi]$) and to satisfy $x > p_x$.

\subsection{Finding intersections}

\subsubsection{Arc–edge intersection (\texttt{isectArcEdge})}

Arc: centre $\mathbf{z}$, radius $\delta$, angles $[\phi_0, \phi_0 + \Delta\varphi]$.
Edge: from $\mathbf{p}$ to $\mathbf{q}$, parametrised as $\mathbf{r}(t) = \mathbf{p} + t\mathbf{e}$,
$\mathbf{e} = \mathbf{q} - \mathbf{p}$, $t \in (0,1)$.

Substituting into the circle equation $|\mathbf{r} - \mathbf{z}|^2 = \delta^2$:
\[
  |\mathbf{e}|^2 t^2 + 2(\mathbf{f}\cdot\mathbf{e})\,t + |\mathbf{f}|^2 - \delta^2 = 0,
  \quad \mathbf{f} = \mathbf{p} - \mathbf{z}.
\]
\[
  a = |\mathbf{e}|^2, \quad b = 2(\mathbf{f}\cdot\mathbf{e}), \quad
  c = |\mathbf{f}|^2 - \delta^2, \quad \Delta = b^2 - 4ac.
\]
Roots $t = (-b \pm \sqrt{\Delta})/(2a)$ (when $\Delta \geq 0$).  Each root with
$t \in (0,1)$ is an intersection; check that the corresponding point lies on the arc.

\subsubsection{Arc–arc intersection (\texttt{isectArcArc})}

Two circles of radii $\delta_A$, $\delta_B$ centred at $\mathbf{z}_A$, $\mathbf{z}_B$.
Let $d = |\mathbf{z}_B - \mathbf{z}_A|$.  The intersection points (up to 2) lie on the
radical axis at distance $a_2 = (d^2 + \delta_A^2 - \delta_B^2)/(2d)$ from $\mathbf{z}_A$,
with perpendicular offset $h = \sqrt{\max(0, \delta_A^2 - a_2^2)}$:
\[
  \mathbf{X} = \mathbf{z}_A + a_2 \hat{\mathbf{s}} \pm h\,\hat{\mathbf{s}}^\perp,
  \quad \hat{\mathbf{s}} = (\mathbf{z}_B - \mathbf{z}_A)/d.
\]
Each candidate is checked against both arc angle ranges.

\subsubsection{Edge–edge intersection (\texttt{isectEdgeEdge})}

Lines $\mathbf{p} + t\mathbf{d}$ and $\mathbf{r} + s\mathbf{e}$:
\[
  t = \frac{(\mathbf{r}-\mathbf{p})\times\mathbf{e}}{\mathbf{d}\times\mathbf{e}},
  \quad
  s = \frac{(\mathbf{r}-\mathbf{p})\times\mathbf{d}}{\mathbf{d}\times\mathbf{e}},
\]
where $\mathbf{a}\times\mathbf{b} = a_x b_y - a_y b_x$.  Intersection if $t, s \in (0,1)$.

%% ============================================================
\section{Force as Gradient of Intersection Area}
%% ============================================================

The force on vertex $\mathbf{a}_\alpha$ is
\[
  \mathbf{F}_\alpha = -\pderiv{E}{\mathbf{a}_\alpha}.
\]
We differentiate $E = \mathrm{Area}(\tilde A \cap \tilde B)$ using the Leibniz
rule for integrals whose limits and integrands both depend on the parameters.

\subsection{General Leibniz rule for signed area}

Consider a smooth arc $\boldsymbol{\gamma}(t;\,\boldsymbol{\theta})$ parameterised by
$t\in[t_0(\boldsymbol{\theta}), t_1(\boldsymbol{\theta})]$ where $\boldsymbol{\theta}$
are the polygon vertex coordinates.  The signed area contribution is
\[
  \mathcal{I}(\boldsymbol{\theta})
  = \frac{1}{2}\int_{t_0}^{t_1}
    \bigl(\gamma_x \dot\gamma_y - \gamma_y \dot\gamma_x\bigr)\,\dd t.
\]
Differentiating with respect to $\boldsymbol{\theta}$:
\begin{align}
  \pderiv{\mathcal{I}}{\boldsymbol{\theta}}
  &= \frac{1}{2}\int_{t_0}^{t_1}
     \pderiv{}{\boldsymbol{\theta}}\bigl(\gamma_x\dot\gamma_y - \gamma_y\dot\gamma_x\bigr)\,\dd t
  \notag\\
  &\quad+ \frac{1}{2}\bigl(\gamma_x(t_1)\dot\gamma_y(t_1)-\gamma_y(t_1)\dot\gamma_x(t_1)\bigr)
     \pderiv{t_1}{\boldsymbol{\theta}}
  \notag\\
  &\quad- \frac{1}{2}\bigl(\gamma_x(t_0)\dot\gamma_y(t_0)-\gamma_y(t_0)\dot\gamma_x(t_0)\bigr)
     \pderiv{t_0}{\boldsymbol{\theta}}.
  \label{eq:leibniz}
\end{align}

The first term is the \textbf{direct gradient} (how the curve itself moves with
$\boldsymbol{\theta}$, holding limits fixed).  The second and third terms are
\textbf{boundary gradient} contributions: the intersection boundary endpoints move
as the polygon vertices change, shifting the limits of integration.

\subsection{Gradient for a straight edge segment}

For the edge section $E_i$ (from $\mathbf{P} = \mathbf{a}_i^+$ to $\mathbf{Q} = \mathbf{a}_{i+1}^-$)
with inside sub-interval $[\mathbf{X}_0, \mathbf{X}_1]$, the energy contribution is
$\mathcal{I} = \frac{1}{2}(X_{0x}X_{1y} - X_{0y}X_{1x})$.

\paragraph{Direct term.}
The segment points are $\boldsymbol{\gamma}(t) = (1-t)\mathbf{P} + t\mathbf{Q}$.
The sub-interval endpoints are $\mathbf{X}_0 = \boldsymbol{\gamma}(t_0)$,
$\mathbf{X}_1 = \boldsymbol{\gamma}(t_1)$.  When $\mathbf{P}$ or $\mathbf{Q}$ moves,
all points on the segment move.  The force on $\mathbf{P}$ from the direct term is:
\[
  \pderiv{\mathcal{I}}{\mathbf{P}}\bigg|_{\text{direct}}
  = \frac{\partial}{\partial\mathbf{P}}\bigl[\tfrac{1}{2}(X_{0x}X_{1y}-X_{0y}X_{1x})\bigr]
    = \frac{1}{2}\bigl[(1-t_0)\,\frac{\partial X_{0x}}{\partial P_x} X_{1y} - \ldots\bigr].
\]
More concisely, define $\mathbf{f}_0 = \frac{1}{2}(X_{1y}, -X_{1x})$ and
$\mathbf{f}_1 = \frac{1}{2}(-X_{0y}, X_{0x})$.  Then:
\begin{align*}
  \pderiv{\mathcal{I}}{\mathbf{P}}\bigg|_{\text{direct}}
    &= (1-t_0)\,\mathbf{f}_0 + (1-t_1)\,\mathbf{f}_1, \\
  \pderiv{\mathcal{I}}{\mathbf{Q}}\bigg|_{\text{direct}}
    &= t_0\,\mathbf{f}_0 + t_1\,\mathbf{f}_1.
\end{align*}
Here the factors $(1-t)$ and $t$ are the barycentric weights of $\mathbf{P}$ and $\mathbf{Q}$
in $\boldsymbol{\gamma}(t)$.

\paragraph{Boundary (limit) term.}
When the endpoints $t_{0,1}$ themselves change (because a crossing parameter
depends on $\mathbf{P}$ or $\mathbf{Q}$), we get an additional contribution.

Consider an edge–edge crossing at parameter $t_0$: the edge of $\tilde A$ (direction
$\mathbf{e}_A = \mathbf{Q}-\mathbf{P}$) crosses an edge of $\tilde B$ (direction $\mathbf{e}_B$)
with $t_0 = (\mathbf{f}\times\mathbf{e}_B)/(\mathbf{e}_A\times\mathbf{e}_B)$.  Then:
\[
  \pderiv{t_0}{\mathbf{P}} = \frac{-\mathbf{e}_B^\perp + t_0\,\mathbf{e}_A^\perp}
                            {\mathbf{e}_A\times\mathbf{e}_B}
  \quad\text{(2-component vector).}
\]
(Here $\mathbf{v}^\perp = (-v_y, v_x)$ is the 90$^\circ$ CCW rotation.)
The Leibniz boundary contribution from $t_0$ is
\[
  \pderiv{\mathcal{I}}{t_0}\cdot\pderiv{t_0}{\mathbf{P}}
  = \frac{1}{2}(\mathbf{e}_A\times\mathbf{X}_1)\cdot\pderiv{t_0}{\mathbf{P}},
\]
and similarly for $t_1$.  In the code these are the ``$\beta$-correction'' terms.

\subsection{Gradient for an arc segment}

For arc $K_i$ with inside sub-interval $[\tau_0, \tau_1]$, the energy contribution is
\[
  \mathcal{I} = \frac{1}{2}(X_{0x}X_{1y} - X_{0y}X_{1x})
              + \frac{1}{2}\delta^2(\Delta\psi - \sin\Delta\psi),
  \quad \Delta\psi = (\tau_1-\tau_0)\Delta\varphi_i.
\]
The arc itself $\boldsymbol{\gamma}(\tau) = \mathbf{z}_i + \delta(\cos\phi(\tau),\sin\phi(\tau))$
moves when the backbone vertices move (because $\mathbf{z}_i$, $\phi_{0,i}$, and $\Delta\varphi_i$
all depend on three vertices: $\mathbf{a}_{i-1}$, $\mathbf{a}_i$, $\mathbf{a}_{i+1}$).
Only the arc \emph{endpoints} $\tau = 0$ ($\mathbf{a}_i^-$) and $\tau = 1$ ($\mathbf{a}_i^+$)
generate simple direct-term gradients (for interior crossings the arc is circular and
a separate chain rule through \texttt{arcGeomGrad} is needed).

\subsection{Arc geometry chain rule (\texttt{arcGeomGrad})}
\label{sec:arcgeomgrad}

The function \texttt{arcGeomGrad} computes, given:
\begin{itemize}
  \item $\pderiv{E}{\mathbf{a}_i^-}$ (force on the tangent point from the incoming edge),
  \item $\pderiv{E}{\mathbf{a}_i^+}$ (force on the tangent point from the outgoing edge),
  \item $\pderiv{E}{\Delta\varphi_i}$ (sensitivity to the arc span),
\end{itemize}
the resulting forces on backbone vertices $\mathbf{a}_{i-1}$, $\mathbf{a}_i$, $\mathbf{a}_{i+1}$.

The chain rule:
\begin{align*}
  \pderiv{E}{\mathbf{a}_{i-1}} &= \pderiv{E}{\mathbf{a}_i^-}\cdot\pderiv{\mathbf{a}_i^-}{\mathbf{a}_{i-1}}
    + \pderiv{E}{\Delta\varphi_i}\cdot\pderiv{\Delta\varphi_i}{\mathbf{a}_{i-1}}, \\
  \pderiv{E}{\mathbf{a}_i} &= \pderiv{E}{\mathbf{a}_i^-}\cdot\pderiv{\mathbf{a}_i^-}{\mathbf{a}_i}
    + \pderiv{E}{\mathbf{a}_i^+}\cdot\pderiv{\mathbf{a}_i^+}{\mathbf{a}_i}
    + \pderiv{E}{\Delta\varphi_i}\cdot\pderiv{\Delta\varphi_i}{\mathbf{a}_i}, \\
  \pderiv{E}{\mathbf{a}_{i+1}} &= \pderiv{E}{\mathbf{a}_i^+}\cdot\pderiv{\mathbf{a}_i^+}{\mathbf{a}_{i+1}}
    + \pderiv{E}{\Delta\varphi_i}\cdot\pderiv{\Delta\varphi_i}{\mathbf{a}_{i+1}}.
\end{align*}

\paragraph{Partial derivatives of $\ell_i$.}
$\ell_i = \delta|c_i|/(1+d_i)$ with $c_i = \uvec_i\times\vvec_i$, $d_i = \uvec_i\cdot\vvec_i$.

Differentiating $\ell$ with respect to $\hat{\mathbf{u}}$ and $\hat{\mathbf{v}}$:
\[
  \pderiv{\ell}{\hat u_x} = \frac{1}{1+d}\bigl(\sgn(c)\,\delta\,\hat v_y - \ell\,\hat v_x\bigr),
  \quad
  \pderiv{\ell}{\hat u_y} = \frac{1}{1+d}\bigl(-\sgn(c)\,\delta\,\hat v_x - \ell\,\hat v_y\bigr),
\]
and similarly for $\pderiv{\ell}{\hat{\mathbf{v}}}$ (swap $\uvec \leftrightarrow \vvec$ with sign changes).

\paragraph{Partials of $\Delta\varphi_i$.}
$\Delta\varphi_i = \mathrm{atan2}(c_i, d_i)$.  Using $c^2 + d^2 = 1$:
\[
  \pderiv{\Delta\varphi}{\hat u_x} = d\hat v_y - c\hat v_x,
  \quad
  \pderiv{\Delta\varphi}{\hat u_y} = -d\hat v_x - c\hat v_y,
\]
and similarly for $\hat{\mathbf{v}}$.

\paragraph{Partials of $\hat{\mathbf{u}}$ with respect to $\mathbf{a}_{i-1}$ and $\mathbf{a}_i$.}
$\hat{\mathbf{u}} = \mathbf{u}/|\mathbf{u}|$ with $\mathbf{u} = \mathbf{a}_i - \mathbf{a}_{i-1}$.
\[
  \pderiv{\hat u_j}{a_{i,k}} = \frac{\delta_{jk} - \hat u_j \hat u_k}{|\mathbf{u}|},
  \qquad
  \pderiv{\hat u_j}{a_{i-1,k}} = -\frac{\delta_{jk} - \hat u_j \hat u_k}{|\mathbf{u}|}.
\]

Combining, the code computes (after careful algebraic simplification) the four
output force contributions:
\[
  \mathbf{f}_{i-1} = -\pderiv{E}{\mathbf{u}} / |\mathbf{u}|, \qquad
  \mathbf{f}_i = \mathbf{f}^{\mathrm{am}} + \mathbf{f}^{\mathrm{ap}}
    + \pderiv{E}{\mathbf{u}}/|\mathbf{u}| - \pderiv{E}{\mathbf{v}}/|\mathbf{v}|, \qquad
  \mathbf{f}_{i+1} = \pderiv{E}{\mathbf{v}} / |\mathbf{v}|,
\]
where $\pderiv{E}{\mathbf{u}} = \bigl(\pderiv{E}{\hat{\mathbf{u}}} - (\pderiv{E}{\hat{\mathbf{u}}}\cdot\uvec)\uvec\bigr)/|\mathbf{u}|$
is the force projected perpendicular to $\uvec$ (since $|\hat{\mathbf{u}}|=1$).

%% ============================================================
\section{SHAKE Constraint Projection}
%% ============================================================

\subsection{Constraints}

For polygon $p$ with $n$ vertices $\{\mathbf{x}_0,\ldots,\mathbf{x}_{n-1}\}$ there are
$n+1$ holonomic constraints:
\begin{align}
  C_k(\mathbf{x}) &= |\mathbf{x}_{k+1} - \mathbf{x}_k| - L_k = 0, \quad k = 0,\ldots,n-1, \\
  C_n(\mathbf{x}) &= A(\mathbf{x}) - A_0 = 0.
\end{align}
The constraint Jacobian $J \in \R^{(n+1)\times 2n}$ has rows $\nabla_\mathbf{x} C_k$.

\subsection{Newton step}

SHAKE applies one Newton step per iteration to drive all constraints to zero simultaneously.

The Newton step is: find $\boldsymbol{\lambda} \in \R^{n+1}$ solving
\[
  (J J^\top)\,\boldsymbol{\lambda} = \mathbf{C}(\mathbf{x}),
\]
then apply the position update
\[
  \Delta\mathbf{x} = -J^\top\boldsymbol{\lambda}.
\]

\subsection{Jacobian structure}

\paragraph{Edge constraints.}
Let $\hat{\mathbf{e}}_k = (\mathbf{x}_{k+1} - \mathbf{x}_k)/|\mathbf{x}_{k+1}-\mathbf{x}_k|$
be the unit edge vector.  Then:
\[
  \nabla_{\mathbf{x}_k} C_k = -\hat{\mathbf{e}}_k, \qquad
  \nabla_{\mathbf{x}_{k+1}} C_k = +\hat{\mathbf{e}}_k, \qquad
  \nabla_{\mathbf{x}_j} C_k = \mathbf{0} \text{ for } j \neq k, k+1.
\]

\paragraph{Area constraint.}
The shoelace formula gives $\nabla_{\mathbf{x}_k} A$.  Using $\mathbf{x}_0$ as reference and
wrapping all other positions:
\begin{align*}
  \pderiv{A}{x_k} &= \tfrac{1}{2}(y_{k+1} - y_{k-1}), \\
  \pderiv{A}{y_k} &= \tfrac{1}{2}(x_{k-1} - x_{k+1}).
\end{align*}
Denote this gradient vector $\mathbf{g}_k = \nabla_{\mathbf{x}_k} A$.

\subsection{Gram matrix $JJ^\top$}

The $(n+1)\times(n+1)$ matrix $M = JJ^\top$ has entries $M_{ij} = \nabla C_i \cdot \nabla C_j$.

\begin{itemize}
  \item \textbf{Diagonal, edge–edge}: $M_{kk} = |\nabla_\mathbf{x} C_k|^2 = 2$
        (the gradient has one $-\hat{\mathbf{e}}_k$ and one $+\hat{\mathbf{e}}_k$, total squared norm $2$).
  \item \textbf{Off-diagonal, adjacent edges} ($j = k\pm 1\mod n$):
        The only shared vertex between $C_k$ and $C_{k+1}$ is $\mathbf{x}_{k+1}$.
        \[M_{k,k+1} = \nabla C_k|_{\mathbf{x}_{k+1}} \cdot \nabla C_{k+1}|_{\mathbf{x}_{k+1}}
          = \hat{\mathbf{e}}_k \cdot (-\hat{\mathbf{e}}_{k+1}) = -\hat{\mathbf{e}}_k \cdot \hat{\mathbf{e}}_{k+1}.\]
  \item \textbf{Off-diagonal, non-adjacent edges}: $M_{k,j} = 0$ for $|k-j| > 1$ (mod $n$,
        except for cyclic adjacency).
  \item \textbf{Area row and column}: $M_{n,k} = M_{k,n} = \nabla C_k \cdot \nabla C_n$.
        \[M_{n,k} = -\hat{\mathbf{e}}_k\cdot\mathbf{g}_k + \hat{\mathbf{e}}_k\cdot\mathbf{g}_{k+1}
          = \hat{\mathbf{e}}_k\cdot(\mathbf{g}_{k+1}-\mathbf{g}_k).\]
  \item \textbf{Area diagonal}: $M_{n,n} = |\nabla_\mathbf{x} A|^2 = \sum_k |\mathbf{g}_k|^2$.
\end{itemize}

\subsection{LU solve and position update}

The $(n+1)\times(n+1)$ system $M\boldsymbol{\lambda} = \mathbf{C}$ is solved by LU factorization
(no pivoting, as $M$ is typically well-conditioned for equilateral polygons).  Then:
\[
  \Delta x_k = \lambda_k\,\hat e_{k,x} - \lambda_{k-1}\,\hat e_{k-1,x} - \lambda_n g_{k,x},
\]
\[
  \Delta y_k = \lambda_k\,\hat e_{k,y} - \lambda_{k-1}\,\hat e_{k-1,y} - \lambda_n g_{k,y}.
\]
(Thread $k$ handles vertex $k$; all threads run in parallel within one block.)

\subsection{Convergence}

Convergence is checked by $\max_k |C_k(\mathbf{x})| < \mathrm{tol}$ after each
iteration, with $\mathrm{tol} = 10^{-15}$.  Up to \texttt{maxIter} iterations are
applied.  If \texttt{maxIter} is reached, the code retries with $10\times$ more
iterations; if it still fails, the FIRE step is rejected (positions reverted).

%% ============================================================
\section{Force Projection}
%% ============================================================

\subsection{Goal}

After computing the overlap force $\mathbf{F}$, we project out all components that
lie in the \emph{constraint subspace} (the span of the constraint gradients).  This
ensures FIRE moves the system along the constraint manifold, not toward constraint
violations.

The projection onto the orthogonal complement of the constraint subspace is:
\[
  \mathbf{F}_\perp = \mathbf{F} - J^\top (J J^\top)^{-1} J\, \mathbf{F},
\]
where $J$ is the $(n+1)\times 2n$ constraint Jacobian.

\subsection{Edge gradient matrix and SVD}

\paragraph{Build $G$ (\texttt{buildEdgeGradMatrixKernel}).}
For each polygon, construct the $n \times 2n$ matrix $G$ where row $k$ is the
gradient of the $k$-th edge constraint:
\[
  G_{k, 2k}   = -\hat e_{k,x}, \quad G_{k, 2k+1} = -\hat e_{k,y}, \quad
  G_{k, 2(k+1)} = +\hat e_{k,x}, \quad G_{k, 2(k+1)+1} = +\hat e_{k,y},
\]
all other entries zero.

\paragraph{SVD (\texttt{cusolverDnDgesvdaStridedBatched}).}
Compute $G = U \Sigma V^\top$ where $U \in \R^{2n\times n}$ and $V \in \R^{n\times n}$.
The columns $\mathbf{u}_1,\ldots,\mathbf{u}_n$ of $U$ are the left singular vectors
and span the range of $G^\top$, i.e.\ the constraint subspace.

\paragraph{Projection using $U$.}
Because $G^\top(G G^\top)^{-1}G = U U^\top$ (when $G$ has full row rank):
\[
  \mathbf{F}_{\perp,\mathrm{edge}} = \mathbf{F} - U U^\top \mathbf{F}
  = \mathbf{F} - \sum_{j=1}^n (\mathbf{u}_j \cdot \mathbf{F})\,\mathbf{u}_j.
\]

\subsection{Area gradient orthogonalisation (\texttt{gramSchmidtAreaKernel})}

The area gradient $\mathbf{g}_A = \nabla_\mathbf{x} A$ (stored in the \texttt{constraints}
array) is not generally in the span of the edge gradients.  We Gram–Schmidt orthogonalise
$\mathbf{g}_A$ against the columns of $U$:
\[
  \tilde{\mathbf{g}}_A = \mathbf{g}_A - \sum_{j=1}^n (\mathbf{u}_j \cdot \mathbf{g}_A)\,\mathbf{u}_j,
\]
then normalise:
\[
  \hat{\mathbf{q}}_A = \tilde{\mathbf{g}}_A / |\tilde{\mathbf{g}}_A|.
\]
This is the unit vector in the constraint subspace direction corresponding to the area
constraint, orthogonal to all edge constraint directions.

\subsection{Full force projection (\texttt{forceProjectFullKernel})}

Project out both edge and area components:
\[
  \mathbf{F}_\perp = \mathbf{F}
    - \sum_{j=1}^n (\mathbf{u}_j \cdot \mathbf{F})\,\mathbf{u}_j
    - (\hat{\mathbf{q}}_A \cdot \mathbf{F})\,\hat{\mathbf{q}}_A.
\]
The result $\mathbf{F}_\perp$ has no component in the direction of any constraint
gradient, so FIRE will not drive constraint violations.

%% ============================================================
\section{FIRE Minimization}
%% ============================================================

\subsection{Background}

FIRE (Fast Inertial Relaxation Engine, Bitzek \textit{et al.}\ 2006) is a molecular
dynamics–based minimizer that uses an inertial term to accelerate convergence toward
energy minima.  The key idea is to run a modified velocity-Verlet integrator and
periodically align the velocity vector with the force direction.

\subsection{Velocity-Verlet integration}

For a particle at position $\mathbf{x}$ with velocity $\mathbf{v}$ and force $\mathbf{F}$:

\begin{enumerate}
  \item \textbf{Half-velocity update} (\texttt{updatePositionAndVelocityFIREKernel}, first half):
  \[
    \mathbf{v}^{n+1/2} = \mathbf{v}^n + \tfrac{1}{2}\mathbf{F}^n\,\dd t,
  \]
  \item \textbf{Position update} (same kernel):
  \[
    \mathbf{x}^{n+1} = \mathbf{x}^n + \mathbf{v}^{n+1/2}\,\dd t + \tfrac{1}{2}\mathbf{F}^n(\dd t)^2
                     = \mathbf{x}^n + \mathbf{v}^n \dd t + \tfrac{1}{2}\mathbf{F}^n(\dd t)^2,
  \]
  followed by periodic wrap $x \mapsto x - \lfloor x\rfloor$.
  \item \textbf{SHAKE}: project $\mathbf{x}^{n+1}$ onto the constraint manifold.
  \item \textbf{Velocity rederivation} (\texttt{rederiveVelocityFromDisplacementKernel}):
  after SHAKE may move positions significantly, the velocity is rederived from the
  actual displacement:
  \[
    \mathbf{v}^{n+1/2}_{\mathrm{redep}} = \mathrm{wrap}(\mathbf{x}^{n+1}_{\mathrm{SHAKE}} - \mathbf{x}^n) / \dd t.
  \]
  \item \textbf{Force and projection}: compute $\mathbf{F}^{n+1}$ via
  \texttt{updateForceEnergyRounded} + \texttt{projectForce}.
  \item \textbf{Second half-velocity update} (\texttt{updateVelocityFIREKernel}):
  \[
    \mathbf{v}^{n+1} = \mathbf{v}^{n+1/2}_{\mathrm{redep}} + \tfrac{1}{2}\mathbf{F}^{n+1}\,\dd t.
  \]
\end{enumerate}

\subsection{FIRE steering}

After each step, compute the \emph{power}:
\[
  P = \mathbf{F}^{n+1} \cdot \mathbf{v}^{n+1}.
\]

\paragraph{If $P > 0$ (moving toward lower energy):}
\begin{enumerate}
  \item Bend velocity toward force direction:
  \[
    \mathbf{v} \leftarrow (1-\alpha)\,\mathbf{v}
               + \alpha\,\frac{|\mathbf{v}|}{|\mathbf{F}|}\,\mathbf{F}.
  \]
  \item Increment $n_+$ (the count of consecutive positive-power steps).
  If $n_+ \geq n_{\min}$, increase $\dd t$:
  \[
    \dd t \leftarrow \min(\dd t \cdot f_{\mathrm{inc}},\; \dd t_{\max}),
    \qquad \alpha \leftarrow \alpha \cdot f_\alpha.
  \]
\end{enumerate}

\paragraph{If $P \leq 0$ (moving uphill):}
\begin{enumerate}
  \item Zero the velocity: $\mathbf{v} \leftarrow \mathbf{0}$.
  \item Decrease $\dd t$:
  \[
    \dd t \leftarrow \dd t \cdot f_{\mathrm{dec}}.
  \]
  \item Reset $\alpha \leftarrow \alpha_{\mathrm{start}}$, $n_+ \leftarrow 0$.
\end{enumerate}

\paragraph{Parameters (defaults used in the packing test):}
\begin{center}
\begin{tabular}{lll}
  Parameter & Symbol & Value \\
  \hline
  Initial time step & $\dd t_0$ & $10^{-3}$ \\
  Maximum time step & $\dd t_{\max}$ & $5\times10^{-2}$ \\
  Initial velocity mixing & $\alpha_{\mathrm{start}}$ & $0.1$ \\
  Mixing decay & $f_\alpha$ & $0.99$ \\
  $\dd t$ increase factor & $f_{\mathrm{inc}}$ & $1.1$ \\
  $\dd t$ decrease factor & $f_{\mathrm{dec}}$ & $0.5$ \\
  Min steps before increase & $n_{\min}$ & $5$ \\
\end{tabular}
\end{center}

\subsection{Time-step scaling for dense packings}

For a system with maximum force $F_{\max}$ and minimum edge length $L_{\min}$,
a safe time step is determined by requiring that the maximum displacement per step
at $\dd t_{\max}$ remains small compared to $L_{\min}$:
\[
  v_{\max}\,\dd t_{\max} \lesssim \eps\,L_{\min},
  \qquad v_{\max} \sim F_{\max}\,\dd t_{\max}.
\]
This gives
\[
  \dd t_{\max} \lesssim \sqrt{\eps L_{\min} / F_{\max}}.
\]
For $\eps = 10^{-3}$, $L_{\min} \approx \mathrm{EL}_{\mathrm{small}}$, and
$F_{\max}$ from the initial configuration.

%% ============================================================
\section{Shape Generation via SLSQP}
%% ============================================================

\subsection{Goal}

Generate a random equilateral closed polygon with $N$ sides, shape index $\kappa$,
and backbone area $A_{\mathrm{tgt}}$, by optimizing a random distribution of interior
angles subject to closure, equilateral, and area constraints.

\subsection{Angle parameterisation}

Parameterise the polygon by its $N-1$ free interior turning angles
$\boldsymbol{\phi} = (\phi_0, \ldots, \phi_{N-2}) \in (0, \pi)^{N-1}$ (the $N$-th angle
is fixed by the closure condition $\sum_k \phi_k = 2\pi$ minus the last).

Given $\boldsymbol{\phi}$, construct the vertices:
\[
  \theta_k = \sum_{j=0}^{k-1} \phi_j, \quad
  \hat{\mathbf{e}}_k = (\cos\theta_k, \sin\theta_k), \quad
  \mathbf{v}_0 = \mathbf{0}, \quad
  \mathbf{v}_{k+1} = \mathbf{v}_k + \hat{\mathbf{e}}_k.
\]
(Unit edge lengths; all $L_k = 1$.)

\subsection{Closure constraint}

For the polygon to close: $\mathbf{v}_N = \mathbf{v}_0 = \mathbf{0}$, i.e.:
\[
  \sum_{k=0}^{N-1} \hat{\mathbf{e}}_k = \mathbf{0}.
\]
This gives 2 equality constraints on $\boldsymbol{\phi}$.

\subsection{Shape index constraint}

The polygon has all edges of unit length, so perimeter $P = N$.  The area is
$A = A(\boldsymbol{\phi})$.  The shape index is $\kappa = P/\sqrt{A} = N/\sqrt{A}$.
Requiring $\kappa = \kappa_0$ is equivalent to requiring
\[
  A(\boldsymbol{\phi}) = (N/\kappa_0)^2 =: A_{\mathrm{unit}}.
\]
This is 1 additional equality constraint.

\subsection{SLSQP formulation}

\begin{align*}
  \min_{\boldsymbol{\phi}} \quad & \tfrac{1}{2}|\boldsymbol{\phi} - \boldsymbol{\phi}_0|^2
    \quad\text{(stay close to random reference)} \\
  \text{s.t.} \quad
  & \sum_k \cos\theta_k(\boldsymbol{\phi}) = 0, \\
  & \sum_k \sin\theta_k(\boldsymbol{\phi}) = 0, \\
  & A(\boldsymbol{\phi}) = A_{\mathrm{unit}}, \\
  & \pi > \sum_k \phi_k > \pi, \quad \phi_k \in (10^{-4}, \pi - 10^{-4}).
\end{align*}
(The range constraint on $\sum \phi_k$ is a convenience to avoid self-intersections
and ensure a convex-ish polygon; the bounds on $\phi_k$ prevent degenerate angles.)

The reference $\boldsymbol{\phi}_0$ is drawn from a Dirichlet distribution scaled to
$[0, 2\pi]$, giving random convex-ish shapes.  If SLSQP succeeds, the resulting
unit-area polygon is scaled to area $A_{\mathrm{tgt}}$:
\[
  \mathbf{v}_k \leftarrow \sqrt{A_{\mathrm{tgt}}/A_{\mathrm{unit}}}\,(\mathbf{v}_k - \bar{\mathbf{v}}),
\]
where $\bar{\mathbf{v}}$ is the centre of mass.  A random rotation is applied to
break any orientation bias.

%% ============================================================
\section{Bidisperse Packing Setup}
%% ============================================================

\subsection{Target areas}

The packing has $N_L = N_S = N_{\mathrm{poly}}/2$ large and small polygons in a unit
$[0,1)^2$ box.  The \emph{backbone packing fraction} is:
\[
  \phi = N_L A_L + N_S A_S.
\]
The large and small polygons have the same shape index $\kappa$ but the large polygon's
perimeter is $\mathrm{ratio}$ times larger than the small polygon's:
\[
  P_L = \mathrm{ratio}\cdot P_S \implies A_L = \mathrm{ratio}^2 \cdot A_S.
\]
Substituting into the packing fraction:
\[
  \phi = N_L \mathrm{ratio}^2 A_S + N_S A_S = A_S(N_L\,\mathrm{ratio}^2 + N_S),
\]
\[
  \boxed{A_S = \frac{\phi}{N_L\,\mathrm{ratio}^2 + N_S}, \qquad A_L = \mathrm{ratio}^2\,A_S.}
\]

\subsection{Target edge lengths}

For a polygon with shape index $\kappa$, backbone area $A$, and $N$ equilateral edges
of length $L$:
\[
  P = N L = \kappa\sqrt{A} \implies L = \frac{\kappa\sqrt{A}}{N}.
\]
SHAKE uses these target edge lengths to maintain equilateral polygons.

\subsection{Rounding radius}

The rounding radius $\delta$ is chosen relative to the \emph{smaller} polygon's edge length:
\[
  \delta = 0.15\,L_S = \frac{0.15\,\kappa\sqrt{A_S}}{N}.
\]
This ensures the validity condition $\delta < L/(2\tan(\pi/N)) \approx L\cdot N/(2\pi)$
(for large $N$, $\tan(\pi/N) \approx \pi/N$) is comfortably satisfied for both polygon sizes.

\subsection{Packing fraction for the 2-polygon test}

For the two-polygon test we use a reference large polygon area of $A_L = 0.8/64$ (same
as the per-polygon area in a 64-polygon packing at $\phi = 0.8$).  The small polygon
is obtained by $A_S = A_L / \mathrm{ratio}^2$.  In this case:
\[
  A_L = \frac{0.8}{64} = 0.01250, \qquad A_S = \frac{0.01250}{1.96} \approx 0.006378.
\]

%% ============================================================
\section{Periodic Boundary Conditions}
%% ============================================================

All positions live in $[0,1)^2$.  The minimum-image convention is:
\[
  \mathrm{wrap}(x) = x - \lfloor x + \tfrac{1}{2}\rfloor.
\]
When computing edge vectors $\mathbf{u} = \mathbf{a}_k - \mathbf{a}_{k-1}$,
the wrap function is applied component-wise to $a_{k,x} - a_{k-1,x}$ and
$a_{k,y} - a_{k-1,y}$.

For computing the interaction between polygon $A$ (where vertex $i$ lives)
and polygon $B$, all of $B$'s vertices are placed in the same image as $B$'s
first vertex, which is placed at the minimum-image position relative to vertex $i$:
\[
  \mathbf{b}_0^{\mathrm{MIC}} = \mathbf{a}_i + \mathrm{wrap}(\mathbf{b}_0 - \mathbf{a}_i),
  \qquad \mathbf{b}_j^{\mathrm{MIC}} = \mathbf{b}_0^{\mathrm{MIC}} + (\mathbf{b}_j - \mathbf{b}_0).
\]
This ensures the entire polygon $B$ is in a single consistent periodic image, preventing
artefacts from vertices on opposite sides of the wrap boundary.

%% ============================================================
\section{Summary of Computational Flow}
%% ============================================================

\begin{enumerate}
  \item \textbf{Initialization}:
    Generate $N_{\mathrm{poly}}$ random equilateral polygons with target areas and
    shape index via SLSQP.  Place on jittered grid.  Initialize neighbour cells.

  \item \textbf{Force/energy evaluation} (\texttt{updateForceEnergyRoundedCUDA}):
    For each vertex $i$ of polygon $A$, in parallel:
    \begin{enumerate}
      \item Compute arc geometry at $i$ and $\mathrm{next}[i]$.
      \item Find all polygons $B \neq A$ in the 3×3 neighbourhood of $i$'s cell.
      \item For arc $K_i$: find crossings with all edges/arcs of $\tilde B$;
            integrate signed area over inside sub-intervals.
      \item For edge $E_i$: similarly.
      \item Accumulate forces on $\{i\!-\!1, i, i\!+\!1, i\!+\!2\}$ via
            chain rule through arc geometry; also update forces on $B$'s vertices.
    \end{enumerate}

  \item \textbf{Force projection} (\texttt{projectForceCUDA}):
    SVD of edge-gradient matrix + Gram–Schmidt area projection.

  \item \textbf{FIRE step}:
    Half-velocity update $\to$ position update $\to$ SHAKE $\to$ velocity rederivation
    $\to$ force evaluation $\to$ second half-velocity $\to$ FIRE steering.

  \item Repeat until $|\mathbf{F}_\perp|_\infty < F_{\mathrm{tol}}$ or maximum steps.
\end{enumerate}

\end{document}
